% =============================================================
%  ch21_kalman_knoten.tex  --  Kalman-Filter-Knoten (Ballverfolgung)
% =============================================================

\chapter{Kalman-Filter-Knoten (Ballverfolgung)}

\begin{lstlisting}[style=py,caption={EKF/KF-Knoten: 3D-Position rein, geglättete Schätzung + Vorhersage raus}]
import numpy as np, rclpy
from rclpy.node import Node
from geometry_msgs.msg import PointStamped

class BallTracker(Node):
    def __init__(self):
        super().__init__('ball_tracker')
        dt = 1/60.0
        self.x = np.zeros((6,1))              # [x y z vx vy vz]
        self.P = np.eye(6) * 1.0
        self.F = np.eye(6); self.F[:3,3:] = np.eye(3)*dt
        self.g  = np.array([[0],[0],[-9.81]])
        self.Bu = np.vstack([0.5*dt**2*self.g, dt*self.g])
        self.H  = np.zeros((3,6)); self.H[:,:3] = np.eye(3)
        self.Q  = np.eye(6)*0.01
        self.R  = np.eye(3)*0.0025            # ~5 cm Messrauschen
        self.sub   = self.create_subscription(
            PointStamped, '/ball/position', self.update, 10)
        self.timer = self.create_timer(dt, self.predict)

    def predict(self):
        self.x = self.F @ self.x + self.Bu
        self.P = self.F @ self.P @ self.F.T + self.Q

    def update(self, msg):
        z = np.array([[msg.point.x],[msg.point.y],[msg.point.z]])
        S = self.H @ self.P @ self.H.T + self.R
        K = self.P @ self.H.T @ np.linalg.inv(S)   # Kalman-Gain
        self.x = self.x + K @ (z - self.H @ self.x)
        self.P = (np.eye(6) - K @ self.H) @ self.P

    def predict_intercept(self, z_plane):
        # reine Praediktion bis zur Schlagebene (kein Messupdate)
        x = self.x.copy()
        for _ in range(120):
            x = self.F @ x + self.Bu
            if x[2,0] <= z_plane:
                return x[:3].flatten()
        return None
\end{lstlisting}
